\section{Project Background}
\label{sec:1-1}

Retail businesses rely on repeat purchases from existing customers, and a discount is the most common tool to keep a customer who is drifting away. A discount is not free, however. For a customer who would otherwise have stopped buying, a voucher can recover a sale. For a customer who would have bought anyway, the same voucher only reduces the margin on a sale the business would have made without it. Sending vouchers to everyone therefore gives away a large share of the discount budget.

The retailer in this project can send a 10\% voucher to some of its active customers every week, within a fixed discount budget. The value of the decision lies in choosing \emph{who} receives it: customers likely to stop buying are where a voucher can pay off, and customers who buy regularly are where it mostly costs margin. Customer churn prediction estimates, for every active customer and every week, how likely they are to make no purchase in the coming months, which is exactly the information needed to rank customers for the offer.

\section{Problem Statement}
\label{sec:1-2}

\begin{itemize}
    \item \textbf{Who should get a retention discount when the weekly budget is fixed?} Each Monday the retailer must choose whom to contact within a fixed discount budget. A voucher can recover a purchase from a customer who would otherwise stop buying, or subsidise a purchase from a customer who would buy anyway.
    \item \textbf{Churn has to be detected early, without contracts.} Customers never cancel in retail; they simply stop buying. Traditional methods notice inactive customers only after a long silence, when the chance to keep them is already small, so churn risk has to be inferred from purchase behavior before the customer is gone.
    \item \textbf{The value of an offer is not observed.} The data contain no past vouchers, margins or campaign results, so the profit of any targeting rule can only be simulated under stated assumptions and must later be confirmed in a pilot.
\end{itemize}

\section{Research Questions}
\label{sec:1-3}

\begin{enumerate}
    \item How can transaction history be converted into weekly customer-level data for churn prediction?
    \item How well does purchasing behavior predict that a customer makes no purchase in the next 90 days?
    \item At the same expected weekly discount spend, which way of choosing customers adds the most simulated net margin: ranking by churn risk, by expected margin per contact, or by time since the last purchase?
    \item Which customer behaviors are most associated with churn? (Correlation, not causation.)
\end{enumerate}

\section{Project Objectives}
\label{sec:1-4}

To develop a \textbf{profit-driven discount strategy}: a churn model that scores every active customer each week, and a decision rule that uses the score to choose which customers receive a discount within a fixed weekly budget, so as to maximise the expected incremental profit relative to sending no offer. The rule is evaluated on unseen test weeks under explicit assumptions, and the report ends with the pilot needed to measure its real effect.

\section{Scope and Limitations}
\label{sec:1-5}

\textbf{In scope:} existing-customer retention, weekly churn scoring of all active customers, budget-constrained discount targeting, and a scenario simulation of its net margin.

\textbf{Out of scope / not claimed:}
\begin{itemize}
    \item Causal uplift and measured campaign revenue -- the transaction data contain no randomized treatment/control information, so every margin figure is a simulation (Chapter~\ref{ch:evaluation}).
    \item A live deployment. \textbf{Chapter~\ref{ch:targeting-policy} describes a proposed operational design}; it is not evaluated on a running system.
    \item Profit analysis based on real margins or costs -- these are not in the data and enter only as stated assumptions.
\end{itemize}

\section{Report Structure}
\label{sec:1-6}

The report is organized into ten chapters that follow the decision from the business problem to the recommended policy:

\begin{enumerate}
    \item \textbf{Introduction} -- background, problem statement, research questions, objectives and scope.
    \item \textbf{Problem Statement and Approach} (Chapter~\ref{ch:problem-approach}) -- the discount decision, its success criteria, and the analytical approach.
    \item \textbf{Dataset and Data Specification} (Chapter~\ref{ch:dataset}) -- the Online Retail II dataset, its fields, conventions and limitations.
    \item \textbf{Data Understanding} (Chapter~\ref{ch:data-understanding}) -- revenue and customers, the purchase rhythm behind the churn definition, and the customer segments.
    \item \textbf{Data Preparation} (Chapter~\ref{ch:data-preparation}) -- cleaning, weekly populations, churn labels, features and the embargoed temporal split.
    \item \textbf{Modeling} (Chapter~\ref{ch:modeling}) -- model selection, tuning and the evaluation of the final model on the test weeks.
    \item \textbf{Evaluation} (Chapter~\ref{ch:evaluation}) -- the business impact of each way of choosing customers at the same discount budget.
    \item \textbf{Targeting Policy and Decision Support} (Chapter~\ref{ch:targeting-policy}) -- the weekly decision cycle, the score table, the dashboard and monitoring.
    \item \textbf{Robustness, Ethics and Limitations} (Chapter~\ref{ch:robustness}) -- how far the result depends on the assumptions, the season and customer groups.
    \item \textbf{Conclusion and Recommendations} (Chapter~\ref{ch:conclusion}) -- the recommended weekly policy and the pilot that decides whether to scale.
\end{enumerate}
